\section{Cut-and-Project Background: Model Sets, Torus Parametrization, and Symbolic Factors}\label{app:cut_project_background}

This appendix serves solely as a background reference point for the geometric context: standard facts and theorem forms concerning cut-and-project sets (model sets), torus phase parametrization, and Sturmian coding can be found in \cite{BaakeGrimmAperiodicOrder1,MeyerAlgebraicNumbers,QueffelecSubstitution,LindMarcus1995}. The proof chain of the present paper does not depend on any statement herein; this material is included only to align the terminology of ``rotational scan / window readout / stable grammar'' used in the main text with the caliber of the aforementioned literature when needed.

\begin{remark}[Phason parametrization and ``time'' readout (for orientation only)]\label{rem:app-cutproject-phason-time}
In the cut-and-project caliber, the phase can be regarded as a point on the torus \(\mathbb{T}=\RR^{d}/\mathcal{L}\), and the object-layer evolution of the scan phase is represented by a one-parameter translation orbit
\[
\theta(t)=\theta_0+t v\pmod{\mathcal{L}};
\]
see \cite{BaakeGrimmAperiodicOrder1,MeyerAlgebraicNumbers} for the relevant background.
\end{remark}

\subsection{One-Dimensional Rotation and Interval Window Readout (Sturmian Interface)}
Consider the circle rotation \(T(x)=x+\alpha\ (\mathrm{mod}\ 1)\) with interval window \(W=[0,\beta)\); the binary readout \(s_t(x)=\ind_W(T^t x)\) provides the symbolic interface of ``rotation + window'' used in the main text. The associated combinatorial properties, including balance, are standard results; see \cite{LindMarcus1995}.
